% 9_pruebas.tex

\section{Pruebas realizadas y resultados}

La validación del sistema se organizó en dos niveles: primero, verificaciones funcionales del software independientes del hardware; y segundo, un procedimiento de validación por etapas sobre los nodos físicos, previo a la prueba integral del enlace.
Esta separación permite aislar los errores de lógica de los relacionados con el cableado, los buses o el canal inalámbrico.

Cabe aclarar que, durante la implementación, la salida de audio del nodo receptor se realizó con un DAC I2S PCM5102 (\textit{overlay} \texttt{hifiberry-dac}), en lugar del módulo NS4168 planteado en el avance de diseño; el resto de la arquitectura se mantuvo conforme a la propuesta.

\subsection{Verificaciones funcionales sin hardware}

Los módulos \texttt{config.py} y \texttt{common.py} no dependen del hardware, por lo que su comportamiento se verificó directamente en la computadora.
Se comprobaron los siguientes aspectos del procesamiento de audio y del protocolo de aplicación:

\begin{itemize}
    \item Empaquetado y reensamblado: la fragmentación de los datos en paquetes y su posterior concatenación reproduce exactamente la secuencia PCM original y los parámetros enviados en el mensaje \texttt{START} se recuperan correctamente al desempaquetarlo.
    \item Conversión de resolución: la conversión de las muestras de $24$ bits entregadas por el micrófono a $8$ bits por muestra opera según lo previsto.
    \item Reconstrucción WAV: a partir de los bytes PCM recibidos se reconstruye un archivo WAV reproducible mediante el encabezado de $44$ bytes generado con los parámetros conocidos del sistema.
    \item Compresión ADPCM: el códec IMA ADPCM de $4$ bits por muestra ($\sim$\SI{32}{kbps}) reduce aproximadamente a la mitad el número de paquetes respecto al PCM de $8$ bits y el receptor lo detecta automáticamente mediante el campo \texttt{codec} del \texttt{START}.
    \item Tolerancia a pérdida de paquetes: el receptor reproduce el audio incluso cuando faltan algunos bloques, rellenando los huecos con silencio y solo señala error si la pérdida supera el umbral configurado (\texttt{MAX\_LOSS\_PCT}).
\end{itemize}

\subsection{Procedimiento de validación por etapas en hardware}

Una vez instalado el software en cada nodo, se siguió un procedimiento de validación incremental antes de probar el enlace completo:

\begin{enumerate}
    \item Verificación del sistema: se confirma la habilitación del bus SPI, la presencia de la tarjeta de audio I2S (\texttt{arecord -l} en el transmisor, \texttt{aplay -l} en el receptor) y la disponibilidad de las bibliotecas de Python requeridas, después de haber ejecutado el \textit{script} de instalación.
    \item Prueba del micrófono (en transmisor): con el INMP441 cableado, se graba una muestra de prueba para confirmar el alineamiento de bits y ajustar la ganancia de captura.
    \item Ejecución manual de los nodos: se ejecutan \texttt{tx.py} y \texttt{rx.py} de forma directa para observar los mensajes de estado y los posibles errores en pantalla antes de delegar la operación al servicio.
    \item Arranque como servicio: se habilita el servicio \textit{systemd} correspondiente, de modo que cada nodo arranca automáticamente en cada encendido.
\end{enumerate}

\subsection{Pruebas del sistema completo}

La prueba del sistema completo evalúa la captura, transmisión y reproducción de mensajes de voz de aproximadamente \SI{5}{\second}, \SI{10}{\second}  y \SI{15}{\second}, con el códec IMA ADPCM habilitado (que es la opción predeterminada del sistema).
Las tres primeras transmisiones se realizaron a \SI{30}{\meter} con línea de vista, en el pasillo de la Escuela de Ingeniería Eléctrica, donde se contaron 36 pasos y la cuarta a \SI{100}{\meter}, también con línea de vista.
Los parámetros para determinar cada prueba como exitosa, son la correcta reconstrucción del archivo WAV, de forma que se escuche de forma inteligible y que la cantidad de bloques perdidos no supere la tolerancia programada.

La Tabla \ref{tab:tx_resultados} resume los indicadores registrados en el nodo transmisor y la Tabla \ref{tab:rx_resultados} los registrados en el nodo receptor.

\begin{table}[htb]
    \centering
    \caption{Indicadores del nodo transmisor (IMA ADPCM).}
    \label{tab:tx_resultados}
    \begin{tabular}{@{}ccccc@{}}
        \toprule
        Prueba & Distancia & Audio capturado & Paquetes & Resultado \\
        \midrule
        1 & \SI{30}{\meter}  & \SI{20224}{\byte} & 698  & Exitosa \\
        2 & \SI{30}{\meter}  & \SI{41920}{\byte} & 1446 & Exitosa \\
        3 & \SI{30}{\meter}  & \SI{61856}{\byte} & 2133 & Exitosa \\
        4 & \SI{100}{\meter} & \SI{63648}{\byte} & 2195 & Exitosa \\
        \bottomrule
    \end{tabular}
\end{table}

\begin{table}[htb]
    \centering
    \caption{Indicadores del nodo receptor (IMA ADPCM).}
    \label{tab:rx_resultados}
    \begin{tabular}{@{}cccc@{}}
        \toprule
        Prueba & Tiempo TX & Tasa efectiva & Calidad del enlace \\
        \midrule
        1 & \SI{0.79}{\second} & \SI{206.0}{\kilo\bit\per\second} & 698/698 bloques (0\,\%) \\
        2 & \SI{1.70}{\second} & \SI{197.2}{\kilo\bit\per\second} & 1446/1446 bloques (0\,\%) \\
        3 & \SI{2.35}{\second} & \SI{210.1}{\kilo\bit\per\second} & 2133/2133 bloques (0\,\%) \\
        4 & \SI{4.28}{\second} & \SI{119.0}{\kilo\bit\per\second} & 2195/2195 bloques (0\,\%) \\
        \bottomrule
    \end{tabular}
\end{table}

Las tres pruebas para determinar el desempeño mínimo completaron la transferencia y reproducción de audio de forma inteligible con éxito, manteniendo tiempos de transmisión y cantidad de paquetes acordes a la duración de audio capturado en cada prueba.
La tasa efectiva entre \SI{197}{\kilo\bit\per\second} y \SI{210}{\kilo\bit\per\second} se mantuvo por debajo del límite programado de \SI{1}{\mega\bit\per\second} del NRF24L01.
Esta diferencia se debe al ciclo de solicitud-confirmación del \textit{Enhanced ShockBurst}, donde el transmisor debe esperar el \texttt{ACK} de cada bloque antes de enviar el siguiente y la latencia de procesamiento del intérprete de Python entre paquetes impide saturar el canal físico.
En cualquier caso, la tasa efectiva resultó entre $6\times$ y $4\times$ superior a la tasa de audio ADPCM (\SI{32}{kbps}), lo que garantiza que la transferencia se complete en una fracción del tiempo de grabación.

Es importante destacar que el tiempo de transmisión contemplado corresponde al intervalo transcurrido entre la recepción del paquete \texttt{START} y la recepción del paquete \texttt{END} por parte del receptor.
Esto se debe a la condición que debe cumplir el receptor, el cual debe estar en modo de espera y no detecta cuando el transmisor empieza la transmisión, sino hasta que recibe el paquete \texttt{START}.

En este contexto, se obtuvieron tiempos de transmisión cortos para los 3 casos, lo cual fue fundamental para cumplir el requisito de no superar los \SI{150}{\second} entre el inicio de la captura y el final de la reproducción de los $3$ audios en conjunto.
Por otra parte, para estas pruebas no se violó el factor de tolerancia del $5\%$ preestablecido en el código fuente, de forma que no se alcanzó este factor de seguridad, lo cual se refleja en la pérdida nula de bloques.

Para cumplir con los requisitos del desempeño óptimo, se realizó la prueba de los \SI{100}{\meter} con línea de vista, capturando un audio de \SI{15}{\second}, la cual corresponde a la cuarta prueba.
Para ello, se colocó el nodo transmisor en la Facultad de Ciencias Sociales y el nodo receptor se colocó en el Aula Magna.

Los resultados obtenidos demuestran que, a pesar de la gran diferencia de distancia respecto a las pruebas anteriores, la calidad del enlace fue óptima al no haber pérdida de paquetes, gracias a los mecanismos de reenvío de paquetes tanto a nivel del ESB como a nivel de aplicación. 
La prueba a \SI{100}{\meter} confirmó que el enlace mantiene una cobertura y calidad idénticas a las obtenidas a \SI{30}{\meter} bajo condiciones de línea de vista, con la diferencia que el tiempo de transmisión aumentó considerablemente.

En el siguiente \href{https://youtu.be/31yQ3FNfaJI?si=LQ8k5ROzQWWf9_tA}{video demostrativo} se presenta el funcionamiento en tiempo real del sistema completo durante una transmisión real para evidenciar los mecanismos de detección de errores.
En la parte izquierda del monitor se observa el nodo transmisor, mientras que en la parte derecha se aprecia el nodo receptor.
Durante la ejecución, el transmisor captura el mensaje de audio, aplica el proceso de compresión IMA ADPCM y fragmenta la información antes de iniciar la transmisión mediante el enlace de radiofrecuencia.
Para esta prueba, el archivo de audio original corresponde a $65856$ B, los cuales, después del proceso de compresión, se transmiten en un total de $2244$ paquetes.

De forma simultánea, el nodo receptor registra la cantidad de paquetes recibidos, realiza el proceso de descompresión y reconstruye el archivo de audio correspondiente.
Los indicadores mostrados durante la ejecución evidencian que la cantidad de paquetes transmitidos coincide con la cantidad de paquetes recibidos por el receptor.
Asimismo, una vez finalizado el proceso de descompresión, el tamaño del archivo reconstruido corresponde nuevamente a $65856$ B, lo que confirma que la información original fue recuperada y que no se presentaron pérdidas de paquetes durante la transmisión. 
Sin embargo, como se observa en el video, el tiempo de transmisión aumentó de \SI{2.48}{\second} en condiciones de separación cercanas y aumentó a \SI{9.50}{\second} para la prueba con una distancia mayor y obstrucciones.

El resto de demostraciones fueron realizadas durante la prueba presencial del proyecto con supervisión del profesor.